\chapter{Production de chaleur}

% ============================================================
\section{Objectifs pédagogiques}
% ============================================================

\subsection*{Compétences GCF mobilisées}
\begin{itemize}
  \item \textbf{C2-1 à C2-3} — Analyser les systèmes de production de chaleur : identifier composants, chaînes d'énergie, paramètres de fonctionnement
  \item \textbf{C3-2 et C3-3} — Comparer et sélectionner un générateur de chaleur selon critères technico-économiques et environnementaux ; dimensionner la puissance de production
  \item \textbf{C5-1 à C5-3} — Appliquer les réglementations thermiques et environnementales liées aux générateurs
  \item \textbf{C7-4 à C7-6} — Mesurer les performances en mise en service (rendement, COP), interpréter les écarts
  \item \textbf{C8-1 à C8-4} — Vérifier les performances saisonnières, proposer des optimisations
\end{itemize}

\subsection*{Savoirs associés mobilisés}
\begin{itemize}
  \item \textbf{A8} — Combustion appliquée : excès d'air, fumées, condensation, rendement (niveau 3)
  \item \textbf{A5} — Thermodynamique appliquée — machine frigorifique côté chaud (niveau 2–3)
  \item \textbf{B1-1} — Équipements de chauffage — production : chaudières, PAC, réseau urbain, échangeurs (niveau 3)
  \item \textbf{B6-3} — Cogénération / micro-cogénération (niveau 2)
  \item \textbf{S5.5} — Réglementations fluides thermiques : chauffage, combustibles (niveau 3)
  \item \textbf{S10.5 et S10.6} — Mise en service et critères de bon fonctionnement (niveau 3)
\end{itemize}

% ============================================================
\section{Introduction}
% ============================================================

La production de chaleur est le point de départ de tout système CVC. Le technicien supérieur GCF doit être capable de dimensionner la puissance nécessaire, de sélectionner la technologie la plus adaptée (chaudière, PAC, solaire, réseau urbain) et de vérifier les performances en exploitation.

Le choix du générateur engage l'installation pour 15 à 25 ans. Il dépend des besoins thermiques du bâtiment, de l'énergie disponible (gaz, électricité, bois, réseau chaleur), des contraintes réglementaires (RE2020, labels énergétiques) et de la rentabilité économique. En 2024, la RE2020 oriente fortement vers les pompes à chaleur et les énergies renouvelables, au détriment des chaudières à gaz dans les logements neufs.

% ============================================================
\section{Combustion et chaudières à gaz}
% ============================================================

\subsection{Principes de la combustion}

La combustion du gaz naturel (méthane : CH\textsubscript{4}) est une réaction d'oxydation exothermique :

\begin{equation}
  \text{CH}_4 + 2\,\text{O}_2 \rightarrow \text{CO}_2 + 2\,\text{H}_2\text{O} + \text{chaleur}
\end{equation}

\textbf{Volume d'air stœchiométrique} pour brûler 1\,m³ de gaz naturel : environ 9,5\,m³ d'air (valeur standard réseau GDF-Suez).

\begin{definition}[title=Pouvoir calorifique]
\begin{itemize}
  \item \textbf{PCI} (Pouvoir Calorifique Inférieur) : énergie dégagée en supposant que la vapeur d'eau des fumées reste à l'état gazeux. Pour le gaz naturel : PCI $\approx \SI{10{,}55}{\kilo\watt\heure\per\cubic\metre}$.
  \item \textbf{PCS} (Pouvoir Calorifique Supérieur) : énergie dégagée en récupérant également la chaleur de condensation de la vapeur d'eau. PCS $\approx \SI{11{,}7}{\kilo\watt\heure\per\cubic\metre}$.
  \item Rapport : $\text{PCS} / \text{PCI} \approx 1{,}11$ pour le gaz naturel (chaleur latente de condensation de l'eau).
\end{itemize}
\end{definition}

\begin{definition}[title=Excès d'air]
Pour assurer une combustion complète, on fournit plus d'air que la quantité stœchiométrique théorique. L'excès d'air $e$ est :
\begin{equation}
  e = \frac{V_{air\,réel} - V_{air\,stœchio}}{V_{air\,stœchio}} \times 100 \qquad [\%]
\end{equation}
Un excès d'air trop faible → combustion incomplète (CO, suies, risque intoxication). Trop élevé → refroidissement excessif des fumées → pertes par les fumées chaudes → rendement dégradé.

\textbf{Valeurs typiques :} brûleur atmosphérique 20–40\,\%, brûleur à prémélange modulant 5–15\,\%.
\end{definition}

\begin{table}[H]
  \centering
  \caption{Composition des fumées de gaz naturel selon l'excès d'air (valeurs indicatives)}
  \begin{tabular}{lccccc}
    \toprule
    Excès d'air & CO\textsubscript{2} (\%) & O\textsubscript{2} (\%) & CO (ppm) & $T_{point\,rosée}$ & Indice Bacharach \\
    \midrule
    0\,\% (stœchio) & 11,7 & 0 & élevé & 57\,°C & — \\
    10\,\% & 10,7 & 1,6 & traces & 57\,°C & 0 \\
    20\,\% & 9,8 & 3,2 & 0 & 57\,°C & 0 \\
    40\,\% & 8,3 & 6,0 & 0 & 55\,°C & 0 \\
    \bottomrule
  \end{tabular}
\end{table}

\subsection{Rendement des chaudières}

\begin{definition}[title=Rendement de combustion]
\begin{equation}
  \eta_{comb} = 1 - \frac{\text{pertes par les fumées} + \text{pertes par incombus\-tion}}{\text{PCI}} \qquad [\%]
\end{equation}
\end{definition}

\begin{definition}[title=Rendement saisonnier (SCOP thermique)]
Le rendement saisonnier intègre les pertes en veille, les cycles de démarrage et les pertes en eau chaude sanitaire. Il est obligatoirement affiché sur l'étiquette énergie EU 811/2013.
\end{definition}

\begin{table}[H]
  \centering
  \caption{Comparatif des technologies de chaudières gaz}
  \begin{tabular}{p{3.5cm}p{3cm}p{3cm}p{2.5cm}}
    \toprule
    Technologie & Rendement PCI & Principe & RE2020 logement \\
    \midrule
    Chaudière basse température & 88–92\,\% & Eau retour $> 45$\,°C & Interdit neuf \\
    Chaudière à condensation & 95–109\,\% & Récupère chaleur latente fumées & Interdit neuf depuis 2022 \\
    Micro-cogénération gaz & 85–92\,\% therm. & Produit électricité + chaleur & Cas par cas \\
    \bottomrule
  \end{tabular}
\end{table}

\begin{attention}
Depuis le 1\textsuperscript{er} janvier 2022, la RE2020 interdit les chaudières gaz comme système principal de chauffage dans les \textbf{logements neufs}. Elles restent autorisées en rénovation et dans le tertiaire. Le marché de la rénovation reste donc très important pour le technicien GCF.
\end{attention}

\subsection{Chaudières à condensation — fonctionnement}

La chaudière à condensation récupère la chaleur latente de condensation de la vapeur d'eau contenue dans les fumées. Cette condensation se produit lorsque la température des fumées descend en dessous du \textbf{point de rosée} (environ 57\,°C pour le gaz naturel).

\begin{methode}[title=Conditions de fonctionnement en condensation]
La condensation est effective si la température de retour eau est inférieure à 57\,°C. Pour maximiser le rendement :
\begin{itemize}
  \item Température départ/retour : 60/45\,°C ou 50/30\,°C (plancher chauffant)
  \item Modulation de puissance : le brûleur à prémélange module de 20 à 100\,\%
  \item Neutralisation des condensats : pH 3 à 5 → kit de neutralisation obligatoire (Loi sur l'eau)
  \item Évacuation des fumées : conduit plastique PVC ou inox (fumées condensées corrosives)
\end{itemize}
\end{methode}

% ---- Diagramme TikZ : Schéma fonctionnel chaudière à condensation ----
\begin{figure}[H]
  \centering
  \begin{tikzpicture}[
      node distance=1.4cm and 2.0cm,
      bloc/.style={draw, rectangle, rounded corners=3pt, minimum width=2.4cm,
                   minimum height=0.9cm, align=center, font=\small\bfseries,
                   fill=#1!12, draw=#1, thick},
      arr/.style={-{Stealth[length=6pt]}, thick},
      lbl/.style={font=\footnotesize, midway}
    ]

    % Composants chaudière (colonne centrale)
    \node[bloc=FedBlue]  (bruleur) {Brûleur\\prémélange};
    \node[bloc=FedOrange, below=of bruleur] (echangeur) {Échangeur\\primaire};
    \node[bloc=FedGreen,  below=of echangeur] (echange2) {Échangeur\\condenseur};
    \node[bloc=FedRed,    below=of echange2]  (siphon)   {Siphon +\\neutralisateur};

    % Circuits extérieurs
    \node[font=\small, above=0.6cm of bruleur, text=FedBlue] (gaz) {Gaz naturel};
    \node[font=\small, right=2.5cm of echangeur] (dep) {Circuit\\chauffage\\(départ)};
    \node[font=\small, right=2.5cm of echange2]  (ret) {Circuit\\chauffage\\(retour $<57$\,°C)};
    \node[font=\small, below=0.5cm of siphon,  text=FedGreen] (cond) {Condensats neutralisés\\(pH 6,5–7) → évacuation};
    \node[font=\small, left=2.2cm of echange2, text=FedOrange] (fum) {Fumées condensées\\→ conduit PVC/inox};

    % Flèches internes
    \draw[arr, FedBlue]   (gaz)      -- (bruleur);
    \draw[arr, FedOrange] (bruleur)  -- node[lbl, right]{gaz chauds}  (echangeur);
    \draw[arr, FedOrange] (echangeur)-- node[lbl, right]{fumées $\downarrow T$} (echange2);
    \draw[arr, FedGreen]  (echange2) -- node[lbl, right]{condensats}  (siphon);
    \draw[arr, FedGreen]  (siphon)   -- (cond);

    % Circuit eau
    \draw[arr, FedOrange] (echangeur.east) -- (dep);
    \draw[arr, FedBlue]   (ret) -- (echange2.east);

    % Sortie fumées
    \draw[arr, FedOrange, dashed] (echange2.west) -- (fum);

    % Annotation point de rosée
    \draw[FedGreen, dashed, thick] (-1.5,-3.0) -- (3.5,-3.0)
      node[right, font=\footnotesize, text=FedGreen]{Point de rosée $\approx 57$\,°C};
  \end{tikzpicture}
  \caption{Schéma fonctionnel d'une chaudière à condensation (récupération de la chaleur latente des fumées)}
  \label{fig:chaudiere-condensation}
\end{figure}

\subsubsection{Chaudières biomasse (bois)}

\begin{table}[H]
  \centering
  \caption{Comparatif des chaudières bois — caractéristiques}
  \begin{tabular}{p{3cm}p{3.5cm}p{3.5cm}p{2.5cm}}
    \toprule
    Type & Combustible & Rendement & Automatisation \\
    \midrule
    Bûches & Bois bûches (H $\leq 20\,\%$) & 75–85\,\% & Manuelle \\
    Plaquettes & Plaquettes forestières (H 25–35\,\%) & 80–90\,\% & Automatique \\
    Granulés (pellets) & Granulés normés ENplus & 85–95\,\% & Entièrement automatique \\
    \bottomrule
  \end{tabular}
\end{table}

% ============================================================
\section{Pompes à chaleur (PAC)}
% ============================================================

\subsection{Principe thermodynamique}

Une pompe à chaleur transfère de la chaleur d'une source froide (air extérieur, sol, nappe phréatique) vers une source chaude (circuit de chauffage), en utilisant un apport d'énergie électrique.

\begin{definition}[title=Coefficient de performance (COP)]
\begin{equation}
  COP = \frac{\dot{Q}_{chaud}}{P_{élec}} = \frac{\dot{Q}_{froid} + P_{élec}}{P_{élec}}
  \label{eq:cop}
\end{equation}
\begin{itemize}
  \item $\dot{Q}_{chaud}$ : puissance thermique fournie au circuit de chauffage (\si{\kilo\watt})
  \item $P_{élec}$ : puissance électrique consommée (\si{\kilo\watt})
  \item $\dot{Q}_{froid}$ : puissance prélevée sur la source froide (\si{\kilo\watt})
\end{itemize}
Un COP = 3 signifie que pour \SI{1}{\kilo\watt} d'électricité consommée, la PAC produit \SI{3}{\kilo\watt} de chaleur.
\end{definition}

\begin{definition}[title=COP de Carnot (COP maximal théorique)]
\begin{equation}
  COP_{Carnot} = \frac{T_{chaud}}{T_{chaud} - T_{froid}} \qquad \text{(températures en Kelvin)}
  \label{eq:cop_carnot}
\end{equation}
Le COP réel est toujours inférieur au COP de Carnot (irréversibilités). Rapport typique : $COP_{réel} / COP_{Carnot} \approx 0{,}4$ à $0{,}6$.
\end{definition}

\subsection{Cycle frigorifique côté chaud}

Le cycle à compression de vapeur comporte 4 étapes :

\begin{enumerate}
  \item \textbf{Évaporation (source froide)} : le fluide frigorigène s'évapore en prélevant de la chaleur à la source froide. Pression basse, température basse.
  \item \textbf{Compression} : le compresseur élève la pression et la température du fluide (apport d'énergie électrique).
  \item \textbf{Condensation (source chaude)} : le fluide cède sa chaleur au circuit de chauffage en se condensant. \textbf{C'est ici que la chaleur utile est produite.}
  \item \textbf{Détente} : le fluide se détend dans le détendeur (électronique, thermostatique ou à tube capillaire), sa pression et sa température chutent. Retour à l'état initial.
\end{enumerate}

\begin{table}[H]
  \centering
  \caption{Fluides frigorigènes courants pour PAC résidentielle}
  \begin{tabular}{lcccl}
    \toprule
    Fluide & GWP & Pression basse (bar) & Pression haute (bar) & Statut \\
    \midrule
    R410A & 2088 & 4,9 & 24,3 & En cours d'élimination (F-Gas) \\
    R32 & 675 & 5,8 & 26,1 & Très répandu (PAC air/eau) \\
    R290 (propane) & 3 & 4,7 & 18,5 & GWP très faible, inflammable \\
    R454B & 466 & 5,4 & 24,8 & Successeur R410A \\
    \bottomrule
  \end{tabular}
\end{table}

\subsection{Types de PAC et performances}

\begin{table}[H]
  \centering
  \caption{Comparatif des types de PAC — performances et contextes d'usage}
  \begin{tabular}{p{2.5cm}p{3cm}p{2.5cm}p{4cm}}
    \toprule
    Type & Source froide & COP typique (A7/W35) & Contexte \\
    \midrule
    Air/Air & Air extérieur & 2,5–4,5 & Climatisation réversible \\
    Air/Eau & Air extérieur & 2,8–4,5 & Chauffage + ECS maison/collectif \\
    Eau glycolée/Eau & Sol horizontal & 3,5–5,5 & Maison individuelle, terrain disponible \\
    Eau/Eau & Nappe phréatique & 4,0–6,0 & Performances élevées, contraintes hydrogéol. \\
    PAC sur eaux grises & Eau usées bâtiment & 3,0–4,5 & Immeuble collectif, récupération énergie \\
    \bottomrule
  \end{tabular}
\end{table}

\begin{definition}[title=SCOP (Seasonal COP)]
Le SCOP est le COP moyen calculé sur l'ensemble de la saison de chauffage, selon la norme EN 14825. Il intègre les variations de température extérieure, les cycles de dégivrage (PAC air/eau) et les modes de relève.

\textbf{Valeurs indicatives (zone H2, émetteurs basse température) :}
\begin{itemize}
  \item PAC air/eau haute performance : SCOP 3,5–4,5
  \item PAC géothermique eau glycolée : SCOP 4,0–5,5
\end{itemize}

\textbf{Conditions normalisées d'essai EN 14825 :}
\begin{itemize}
  \item A7/W35 : air à 7\,°C, eau départ à 35\,°C (référence basse température)
  \item A2/W35 : air à 2\,°C, eau départ à 35\,°C (point intermédiaire)
  \item A-7/W35 : air à -7\,°C, eau départ à 35\,°C (froid hivernal)
\end{itemize}
\end{definition}

% ---- Diagramme TikZ : Cycle frigorifique PAC ----
\begin{figure}[H]
  \centering
  \begin{tikzpicture}[
      node distance=2.8cm,
      box/.style={draw, rectangle, rounded corners=4pt, minimum width=2.6cm,
                  minimum height=1.1cm, align=center, font=\small\bfseries,
                  fill=#1!15, draw=#1, thick},
      arr/.style={-{Stealth[length=7pt]}, thick, draw=#1}
    ]

    % Nœuds principaux
    \node[box=FedBlue]   (comp) {Compresseur\\{\footnotesize(élec. $P_{él}$)}};
    \node[box=FedOrange, right=of comp] (cond) {Condenseur\\{\footnotesize(source chaude)}};
    \node[box=FedGreen,  below=of cond] (dete) {Détendeur\\{\footnotesize(valve / cap.)}};
    \node[box=FedBlue,   below=of comp] (evap) {Évaporateur\\{\footnotesize(source froide)}};

    % Flèches fluide frigorigène
    \draw[arr=FedBlue]   (comp) -- node[above, font=\footnotesize]{HP vapeur} (cond);
    \draw[arr=FedOrange] (cond) -- node[right, font=\footnotesize]{HP liquide} (dete);
    \draw[arr=FedGreen]  (dete) -- node[below, font=\footnotesize]{BP liquide} (evap);
    \draw[arr=FedBlue]   (evap) -- node[left,  font=\footnotesize]{BP vapeur}  (comp);

    % Échanges énergétiques
    \draw[arr=FedOrange, dashed] (cond.north) -- ++(0, 0.9)
          node[above, font=\small, text=FedOrange]
          {$\dot{Q}_{chaud} = \dot{Q}_{froid} + P_{él}$};
    \draw[arr=FedBlue, dashed]   (evap.south) -- ++(0,-0.9)
          node[below, font=\small, text=FedBlue]
          {$\dot{Q}_{froid}$ prélevé sur source froide};
    \draw[arr=FedBlue, dashed]   (comp.west) -- ++(-1.1, 0)
          node[left, font=\small, text=FedBlue]{$P_{él}$};

    % Légende COP
    \node[draw=FedOrange, fill=FedOrange!10, rounded corners, font=\small,
          below=0.5cm of evap, xshift=1.4cm, text width=4.8cm, align=center]
      {$\displaystyle COP = \frac{\dot{Q}_{chaud}}{P_{él}} > 1$};
  \end{tikzpicture}
  \caption{Cycle frigorifique d'une pompe à chaleur (cycle à compression de vapeur)}
  \label{fig:cycle-pac}
\end{figure}

\subsection{Point de bivalence}

La puissance des PAC air/eau diminue lorsque la température extérieure baisse. En dessous d'une certaine température (point de bivalence), un générateur d'appoint (résistance électrique, chaudière) prend le relais.

\begin{methode}[title=Détermination du point de bivalence]
\begin{enumerate}
  \item Tracer la courbe de puissance de la PAC en fonction de $T_{ext}$ (fournie par le fabricant)
  \item Tracer la droite des déperditions du bâtiment : $\dot{Q}_{dep} = H_T \cdot (T_{int} - T_{ext})$
  \item L'intersection des deux courbes donne le \textbf{point de bivalence} ($T_{biv}$, $\dot{Q}_{biv}$)
  \item En dessous de $T_{biv}$, l'appoint assure la différence de puissance
\end{enumerate}
\textbf{Valeur courante :} $T_{biv}$ entre $-5$\,°C et $+5$\,°C selon le dimensionnement choisi.

Un point de bivalence bas (−5\,°C) signifie une PAC surdimensionnée qui couvre presque tous les besoins sans appoint : confort maximal, coût d'investissement élevé.
Un point de bivalence haut (+5\,°C) : PAC plus petite, appoint plus sollicité.
\end{methode}

% ============================================================
\section{Solaire thermique}
% ============================================================

\subsection{Principe et composants}

Le solaire thermique capte le rayonnement solaire pour produire de la chaleur, principalement pour l'eau chaude sanitaire (ECS) et en appoint de chauffage.

\begin{table}[H]
  \centering
  \caption{Types de capteurs solaires thermiques}
  \begin{tabular}{p{3cm}p{4cm}p{4.5cm}}
    \toprule
    Type & Principe & Performance / Température \\
    \midrule
    Plan vitré & Absorbeur plan sous vitrage & $\eta_0 \approx 0{,}75$, jusqu'à 80\,°C \\
    Tube sous vide & Tube en verre sous vide, pertes minimales & $\eta_0 \approx 0{,}70$, jusqu'à 120\,°C \\
    Concentrateur & Miroir parabolique, rayonnement focalisé & Températures très élevées, usage industriel \\
    \bottomrule
  \end{tabular}
\end{table}

\begin{definition}[title=Rendement d'un capteur solaire plan]
\begin{equation}
  \eta = \eta_0 - a_1 \cdot \frac{T_m - T_{ext}}{G} - a_2 \cdot \frac{(T_m - T_{ext})^2}{G}
\end{equation}
\begin{itemize}
  \item $\eta_0$ : rendement optique (sans pertes thermiques)
  \item $a_1$, $a_2$ : coefficients de pertes thermiques (fournis par le fabricant, certifiés Solar Keymark)
  \item $T_m$ : température moyenne du fluide dans le capteur (\si{\celsius})
  \item $G$ : irradiance solaire (\si{\watt\per\square\metre}) — en France : 150 à 900\,W/m²
\end{itemize}
\end{definition}

\begin{definition}[title=Taux de couverture solaire]
\begin{equation}
  f_{sol} = \frac{\text{énergie solaire utile annuelle}}{\text{besoins thermiques annuels}} \times 100 \qquad [\%]
\end{equation}
Pour une installation ECS résidentielle bien dimensionnée : $f_{sol} = 50$ à 70\,\%.
\end{definition}

\subsection{Systèmes solaires combinés (SSC)}

Un SSC couvre à la fois les besoins ECS et une partie du chauffage. La surface de capteurs est plus importante (20 à 40\,m² pour une maison). Le taux de couverture annuel global est typiquement de 25 à 40\,\% pour l'ensemble chauffage + ECS.

\begin{table}[H]
  \centering
  \caption{Surfaces de capteurs recommandées selon l'usage (France, zone H2)}
  \begin{tabular}{lccc}
    \toprule
    Usage & Personnes / Surface & Surface capteurs & Volume stockage \\
    \midrule
    ECS seule & 4 personnes & 3–4 m² & 200–300 L \\
    ECS seule & 6 personnes & 5–6 m² & 300–400 L \\
    SSC (ECS + chauffage) & Maison 100 m² RT2012 & 15–20 m² & 1000–2000 L \\
    SSC (ECS + chauffage) & Maison 150 m² RE2020 & 20–30 m² & 2000–3000 L \\
    \bottomrule
  \end{tabular}
\end{table}

% ============================================================
\section{Réseaux de chaleur urbains}
% ============================================================

\begin{definition}[title=Réseau de chaleur urbain]
Un réseau de chaleur distribue de l'eau chaude ou de la vapeur produite par une ou plusieurs chaufferies centralisées (cogénération, géothermie, biomasse, récupération chaleur industrielle) vers des sous-stations d'abonnés via un réseau de canalisations enterrées.

\textbf{Avantages :}
\begin{itemize}
  \item Mutualisation de la production, meilleur rendement global
  \item Facilite l'intégration des énergies renouvelables et de récupération (EnR\&R)
  \item Comptabilisé comme EnR dans la RE2020 si taux EnR\&R $\geq 50\,\%$
\end{itemize}
\end{definition}

\begin{definition}[title=Sous-station d'abonné]
La sous-station est le point de livraison de l'énergie thermique. Elle comprend un échangeur de chaleur qui sépare le réseau primaire (réseau urbain) du réseau secondaire (installation de l'abonné). Le comptage de l'énergie livrée est assuré par un compteur d'énergie thermique (calorimètre).
\end{definition}

% ============================================================
\section{Dimensionnement de la puissance de production}
% ============================================================

\begin{methode}[title=Méthode de calcul de la puissance installée]
\begin{enumerate}
  \item \textbf{Calculer les déperditions thermiques} du bâtiment à la température de base (selon la méthode NF EN 12831)
  \item \textbf{Ajouter les besoins ECS} : $\dot{Q}_{ECS} = \dot{m}_{ECS} \cdot c_p \cdot (T_{ECS} - T_{froide})$
  \item \textbf{Ajouter les besoins de ventilation} (préchauffage de l'air neuf en hiver)
  \item \textbf{Appliquer un coefficient de foisonnement} (0,7 à 0,9 pour le collectif)
  \item \textbf{Puissance installée} = somme des besoins / coefficient foisonnement
\end{enumerate}

\textbf{Règle de majoration courante :} majorer de 10 à 15\,\% pour tenir compte des pertes de distribution et des pointes.
\end{methode}

\begin{table}[H]
  \centering
  \caption{Coefficients de foisonnement ECS selon le type de bâtiment}
  \begin{tabular}{lcc}
    \toprule
    Type de bâtiment & Coefficient de foisonnement & Commentaire \\
    \midrule
    Maison individuelle & 1,0 & Pas de foisonnement \\
    Petit collectif ($\leq$ 10 logements) & 0,9 & Foisonnement faible \\
    Grand collectif (10–50 logements) & 0,7–0,8 & Foisonnement significatif \\
    Tertiaire (bureaux) & 0,6–0,7 & Pointes ECS faibles \\
    Hôtel & 0,5–0,7 & Forte diversité des besoins \\
    \bottomrule
  \end{tabular}
\end{table}

% ============================================================
\section{Exemples résolus}
% ============================================================

\begin{exemple}[title=Rendement d'une chaudière à condensation — PCI et PCS]
\textbf{Énoncé :} Une chaudière à condensation brûle \SI{2{,}5}{\cubic\metre\per\hour} de gaz naturel (PCI = \SI{10{,}55}{\kilo\watt\heure\per\cubic\metre}, PCS = \SI{11{,}7}{\kilo\watt\heure\per\cubic\metre}). Elle délivre une puissance de \SI{24}{\kilo\watt} au circuit de chauffage. Calculer le rendement PCI et le rendement PCS.

\textbf{Puissance chimique PCI :}
\[
  \dot{Q}_{PCI} = 2{,}5 \times 10{,}55 = \SI{26{,}375}{\kilo\watt}
\]

\textbf{Rendement PCI :}
\[
  \eta_{PCI} = \frac{24}{26{,}375} = 0{,}910 = \mathbf{91{,}0\,\%}
\]

\textbf{Puissance chimique PCS :}
\[
  \dot{Q}_{PCS} = 2{,}5 \times 11{,}7 = \SI{29{,}25}{\kilo\watt}
\]

\textbf{Rendement PCS :}
\[
  \eta_{PCS} = \frac{24}{29{,}25} = 0{,}820 = \mathbf{82{,}0\,\%}
\]

\textbf{Interprétation :} Le rendement PCS est inférieur au rendement PCI car le PCS intègre la chaleur latente de condensation (la chaudière la récupère mais le calcul PCS comptabilise plus d'énergie en entrée). Pour comparer les chaudières entre elles en Europe, on utilise le rendement PCS (directive Rendement). Un rendement PCS > 100\,\% est possible pour les chaudières à condensation car la récupération de la chaleur latente est comptabilisée en sortie.
\end{exemple}

\begin{exemple}[title=COP d'une PAC air/eau et comparaison économique]
\textbf{Énoncé :} Une PAC air/eau fonctionne avec : $T_{ext} = \SI{7}{\celsius}$ (A7), $T_{dep} = \SI{35}{\celsius}$ (W35), COP fabricant = 4,1. Puissance électrique absorbée : \SI{4}{\kilo\watt}.

\textit{1 — Puissance thermique produite :}
\[
  \dot{Q}_{chaud} = COP \times P_{élec} = 4{,}1 \times 4 = \SI{16{,}4}{\kilo\watt}
\]
\[
  \dot{Q}_{froid} = \dot{Q}_{chaud} - P_{élec} = 16{,}4 - 4 = \SI{12{,}4}{\kilo\watt} \text{ (prélevé sur l'air ext.)}
\]

\textit{2 — COP de Carnot et rendement exergétique :}
\[
  T_{chaud} = 35 + 273 = \SI{308}{\kelvin}, \quad T_{froid} = 7 + 273 = \SI{280}{\kelvin}
\]
\[
  COP_{Carnot} = \frac{308}{308 - 280} = \frac{308}{28} = 11{,}0
\]
\[
  \eta_{ex} = \frac{COP_{réel}}{COP_{Carnot}} = \frac{4{,}1}{11{,}0} = 0{,}373 = \mathbf{37{,}3\,\%}
\]

\textit{3 — Coût de production par kWh chaleur (prix gaz 0,12\,€/kWh, élec 0,22\,€/kWh) :}
\begin{itemize}
  \item \textbf{PAC :} $c_{PAC} = \frac{0{,}22}{4{,}1} = \SI{0{,}054}{\euro\per\kilo\watt\heure}_{th}$
  \item \textbf{Chaudière gaz ($\eta_{PCI} = 95\,\%$) :} $c_{gaz} = \frac{0{,}12}{0{,}95} = \SI{0{,}126}{\euro\per\kilo\watt\heure}_{th}$
\end{itemize}
La PAC est \textbf{2,3 fois moins chère} à l'usage dans ces conditions. L'écart se réduit si le prix de l'électricité augmente plus vite que le gaz.
\end{exemple}

\begin{exemple}[title=Dimensionnement de la puissance d'une chaudière collective]
\textbf{Énoncé :} Un immeuble collectif de 20 logements de 70\,m² chacun est situé en zone climatique H2 ($T_{base} = -7$\,°C, $T_{int} = 19$\,°C). Le coefficient de déperdition moyen est $G = \SI{1{,}2}{\watt\per\square\metre\per\kelvin}$. Les besoins ECS représentent \SI{8}{\kilo\watt} en pointe. Le coefficient de foisonnement est 0,8.

\textit{1 — Déperditions totales :}
\[
  S_{totale} = 20 \times 70 = \SI{1400}{\square\metre}
\]
\[
  \dot{Q}_{dep} = G \times S \times (T_{int} - T_{base}) = 1{,}2 \times 1400 \times (19 - (-7)) = \SI{43\,680}{\watt} \approx \SI{43{,}7}{\kilo\watt}
\]

\textit{2 — Puissance nécessaire (chauffage + ECS, avec foisonnement) :}
\[
  \dot{Q}_{total} = \frac{\dot{Q}_{dep} + \dot{Q}_{ECS}}{foisonnement} = \frac{43{,}7 + 8}{0{,}8} = \frac{51{,}7}{0{,}8} = \SI{64{,}6}{\kilo\watt}
\]

\textit{3 — Choix :} Le modèle A (\SI{60}{\kilo\watt}) est insuffisant (60 < 64,6). Le \textbf{modèle B (\SI{80}{\kilo\watt})} est retenu — marge de 24\,\%, suffisante sans surdimensionnement excessif.
\end{exemple}

\begin{exemple}[title=Analyse des performances d'une PAC en exploitation]
\textbf{Énoncé :} Une PAC air/eau est mesurée lors de sa mise en service :
\begin{itemize}
  \item Puissance électrique mesurée : \SI{3{,}8}{\kilo\watt}
  \item Température départ eau : \SI{45}{\celsius}, retour eau : \SI{40}{\celsius}
  \item Débit eau : \SI{1200}{\litre\per\hour}, $\rho_{eau} = \SI{992}{\kilogram\per\cubic\metre}$, $c_p = \SI{4179}{\joule\per\kilogram\per\kelvin}$
  \item Température extérieure : \SI{2}{\celsius}
\end{itemize}
Le fabricant annonce COP = 3,8 en A2/W45.

\textit{1 — Puissance thermique produite :}
\[
  \dot{m} = \frac{1200 \times 992}{3600 \times 1000} = \SI{0{,}331}{\kilogram\per\second}
\]
\[
  \dot{Q} = 0{,}331 \times 4179 \times 5 = \SI{6{,}92}{\kilo\watt}
\]

\textit{2 — COP mesuré :}
\[
  COP_{mes} = \frac{6{,}92}{3{,}8} = \mathbf{1{,}82}
\]

\textit{3 — Analyse :} Le COP mesuré (1,82) est très inférieur au COP annoncé (3,8 en A2/W45). Pistes de diagnostic :
\begin{itemize}
  \item Vérifier l'état du circuit frigorifique (charge en fluide frigorigène — sous-charge typique)
  \item Vérifier l'encrassement ou le givrage de l'échangeur air/fluide
  \item Vérifier si la PAC est en mode dégivrage lors de la mesure
  \item Vérifier le débit d'air sur l'évaporateur (filtre encrassé ou obstrué)
  \item Vérifier l'absence de recyclage d'air chaud vers l'aspiration
\end{itemize}
\end{exemple}

% ---- Diagramme TikZ : Comparaison COP PAC vs rendement chaudière ----
\begin{figure}[H]
  \centering
  \begin{tikzpicture}[x=1.6cm, y=0.9cm]
    % Axes
    \draw[-{Stealth}, thick] (0,0) -- (6.5,0)
      node[right, font=\small]{$T_{ext}$ (\si{\celsius})};
    \draw[-{Stealth}, thick] (0,0) -- (0,6.0)
      node[above, font=\small]{Performance (\%)};

    % Graduations axe x : -10, -5, 0, 5, 7, 10, 15
    \foreach \x/\lbl in {0.5/-10, 1.5/-5, 2.5/0, 3.5/5, 4.0/7, 4.5/10, 5.5/15}{
      \draw (\x,0.05) -- (\x,-0.05);
      \node[below, font=\footnotesize] at (\x,0) {\lbl};
    }

    % Graduations axe y : 100, 200, 300, 400, 500
    \foreach \y/\lbl in {1/100, 2/200, 3/300, 4/400, 5/500}{
      \draw (0.05,\y) -- (-0.05,\y);
      \node[left, font=\footnotesize] at (0,\y) {\lbl\,\%};
    }

    % Courbe PAC air/eau — COP×100 (linearly increases with T_ext)
    % T_ext: -10→COP 2.2, -5→2.6, 0→3.1, 5→3.6, 7→4.1, 10→4.4, 15→5.0
    \draw[FedBlue, very thick]
      (0.5,2.2) -- (1.5,2.6) -- (2.5,3.1) -- (3.5,3.6) -- (4.0,4.1) -- (4.5,4.4) -- (5.5,5.0);
    \node[FedBlue, font=\small\bfseries, right] at (5.5,5.0) {PAC air/eau (COP×100)};

    % Ligne chaudière condensation ≈ 105 % PCI = constant
    \draw[FedOrange, very thick, dashed] (0.5,1.05) -- (5.5,1.05);
    \node[FedOrange, font=\small\bfseries, right] at (5.5,1.05)
      {Chaudière condensation ($\eta_{PCI}\approx105\,\%$)};

    % Ligne chaudière basse température ≈ 90 %
    \draw[FedRed, thick, dotted] (0.5,0.90) -- (5.5,0.90);
    \node[FedRed, font=\small, right] at (5.5,0.90)
      {Chaudière basse T° ($\eta\approx90\,\%$)};

    % Annotation point de référence A7/W35
    \draw[FedBlue, dashed] (4.0,0) -- (4.0,4.1);
    \node[FedBlue, font=\footnotesize, above] at (4.0,4.1) {A7/W35};

    % Titre
    \node[font=\small\itshape, above] at (3.0,5.8)
      {Performances comparées en fonction de la température extérieure};
  \end{tikzpicture}
  \caption{Comparaison COP d'une PAC air/eau et rendement d'une chaudière en fonction de la température extérieure}
  \label{fig:comparaison-cop-rendement}
\end{figure}

% ============================================================
\section{Éléments de réglementation}
% ============================================================

\begin{description}
  \item[RE2020 (décret du 29 juillet 2021)] Interdit les chaudières à gaz comme système principal dans les logements neufs depuis le 1\textsuperscript{er} janvier 2022. Favorise les PAC (coefficient $b_{EF} = 2{,}3$ pour l'électricité) et les réseaux de chaleur à EnR\&R. Fixe un seuil d'émissions de GES (Ic\textsubscript{énergie}).

  \item[NF EN 12831] Méthode de calcul des déperditions thermiques de bâtiment — base du dimensionnement des générateurs de chaleur.

  \item[NF EN 14825] Méthode de calcul du SCOP des pompes à chaleur — conditions d'essai normalisées (A7/W35, A2/W35, etc.) pour la comparaison des équipements.

  \item[Règlement EU 813/2013 (écoconception chaudières)] Impose un rendement énergétique saisonnier minimal pour les chaudières de chauffage (classe ErP). Depuis 2015, seules les chaudières à condensation ou haute performance sont autorisées dans l'UE.

  \item[Règlement EU 811/2013 (étiquetage énergie)] Étiquette énergie de A+++ à G pour les générateurs de chaleur — obligatoire sur tout équipement vendu.

  \item[Arrêté du 23 juin 1978 (modifié)] Installations fixes destinées au chauffage et à l'alimentation en eau chaude sanitaire. Règles de sécurité : pression max, température max, soupapes, vases d'expansion.

  \item[DTU 61.1 — NF P45-204] Installations de gaz dans les bâtiments — règles de sécurité, raccordements, ventilation des locaux contenant des appareils à gaz.

  \item[Règlement (UE) 517/2014 (F-Gas)] Réglementation sur les gaz fluorés à fort effet de serre. Interdit progressivement les fluides frigorigènes à GWP élevé (R410A interdit dans les nouveaux équipements à partir de 2025). Impose une attestation de capacité (certificat fluides frigorigènes) pour la manipulation.

  \item[MaPrimeRénov' et CEE] Dispositifs d'aide à la rénovation énergétique : bonus pour les PAC (SCOP $\geq 3{,}5$), bonus pour les chaudières biomasse. Les CEE (Certificats d'Économies d'Énergie) financent la rénovation via les obligations des fournisseurs d'énergie.
\end{description}

\begin{tcolorbox}[
  enhanced, colback=FedBlue!3, colframe=FedBlue!50,
  fonttitle=\bfseries\sffamily, coltitle=FedBlue,
  title={FICHE MÉMO — Production de chaleur}, boxrule=0.8pt, arc=2pt,
  width=\textwidth
]
\textbf{Combustion — Chaudières gaz}
\[
  e = \frac{V_{air\,réel} - V_{air\,stœchio}}{V_{air\,stœchio}} \times 100 \quad [\%]
  \qquad \text{excès d'air (brûleur atm. : 20–40\,\%, prémélange : 5–15\,\%)}
\]
\[
  \eta_{comb} = \frac{\dot{Q}_{utile}}{\dot{Q}_{PCI}} \quad [\%]
  \qquad \text{condensation effective si } T_{retour} < 57\,\text{°C (point de rosée gaz)}
\]

\textbf{Pompe à chaleur (PAC) — Cycle à compression}
\[
  COP = \frac{\dot{Q}_{chaud}}{P_{élec}} \qquad \text{avec } \dot{Q}_{chaud} = \dot{Q}_{froid} + P_{élec} \quad [\si{\kilo\watt}]
\]
\[
  COP_{Carnot} = \frac{T_{chaud}}{T_{chaud} - T_{froid}} \quad \text{(températures en \si{\kelvin})}
  \qquad COP_{réel} \approx 0{,}4 \text{ à } 0{,}6 \times COP_{Carnot}
\]
\[
  SCOP = \frac{E_{thermique\,annuelle}}{E_{électrique\,annuelle}}
  \qquad \text{(NF EN 14825 — intègre variations saisonnières et dégivrage)}
\]

\textbf{Solaire thermique}
\[
  \eta = \eta_0 - a_1 \cdot \frac{T_m - T_{ext}}{G} - a_2 \cdot \frac{(T_m - T_{ext})^2}{G}
  \qquad \text{avec } G = \text{irradiance} [\si{\watt\per\square\metre}]
\]
\[
  f_{sol} = \frac{E_{solaire\,utile}}{E_{besoin\,total}} \times 100 \quad [\%]
  \qquad \text{ECS résidentielle bien dimensionnée : } f_{sol} = 50 \text{ à } 70\,\%
\]

\textbf{Dimensionnement de la puissance}
\[
  \dot{Q}_{ECS} = \dot{m}_{ECS} \cdot c_p \cdot (T_{ECS} - T_{froide}) \quad [\si{\kilo\watt}]
  \qquad \text{déperditions : } \dot{Q}_{dep} = H_T \cdot (T_{int} - T_{ext})
\]
\end{tcolorbox}

\finchapitre
